\chapter{100 câu cục súc về học và tỉnh ngộ}
\markboth{100 câu cục súc về học và tỉnh ngộ}{100 câu cục súc về học và tỉnh ngộ}

\begin{truyen}{Cách đọc chương này}{How to read this chapter}
\chuhoa{M}{ỗi câu ở đây ngắn, cứng và hơi gắt. Đọc để chạm, không phải đọc để vuốt. Nếu một câu làm người đọc khó chịu, rất có thể nó vừa chạm đúng chỗ yếu.}
\textit{(These lines are short, blunt, and slightly harsh. They are meant to hit, not to flatter. If a line irritates you, it may have touched a weak point.)}
\end{truyen}

\section{25 câu về chuyện học không màu hồng}
\begin{enumerate}[leftmargin=*]
\item Học không chắc đổi đời ngay, nhưng dốt thì thường khổ sớm.
\item Đừng hỏi học để làm gì khi em còn chưa biết không học thì sẽ trôi tới đâu.
\item Sợ học khó là bình thường, sợ khó nên bỏ học mới đáng lo.
\item Có những đứa chê học vô dụng vì nó chưa từng dùng hết đầu mình.
\item Học không làm em sang ngay, nhưng nó bớt làm em bị dắt như bò.
\item Ghét học không đáng xấu hổ, ghét học mà còn tự hào mới mệt.
\item Nhiều người chê sách chỉ vì họ lười ngồi yên với chính mình.
\item Học không phải để hơn người, mà để đỡ bị đời nắm gáy.
\item Đừng tưởng cái gì kiếm ra tiền nhanh cũng đáng để bỏ công học theo.
\item Người nói học vô nghĩa thường đang sống nhờ công sức của người từng học nghiêm túc.
\item Càng chưa có gì trong tay càng nên học đàng hoàng.
\item Đầu óc rỗng rất thích gọi sự nỗ lực của người khác là làm màu.
\item Học chậm vẫn đỡ hơn sống nhanh mà nghĩ nông.
\item Nhiều đứa lười đọc vì sợ trong đầu mình không có gì để đối thoại.
\item Học không phải cứu em khỏi mọi khổ, nhưng bớt cho em khổ ngu.
\item Có những cái giá chỉ người thiếu học mới phải trả nhiều lần.
\item Không thích học vẫn phải học, vì đời không quan tâm em có thích hay không.
\item Kỹ năng không tự mọc ra chỉ vì em lớn thêm tuổi.
\item Học là lúc em còn được sai rẻ, ra đời sai thường đắt hơn.
\item Càng sợ bị coi thường càng phải tự bồi đầu mình lên.
\item Đừng lấy vài người giàu không học làm cớ để tự cho phép mình rỗng.
\item Nhiều người bảo thực tế rồi khuyên em bỏ học, thật ra họ chỉ đang quen sống ngắn hạn.
\item Học không làm em hết lạc, nhưng cho em bản đồ để tự tìm đường.
\item Sợ học là chuyện nhỏ, sợ lớn mà vẫn rỗng mới là chuyện lớn.
\item Muốn đời đỡ vả thì đừng ghét cái đang giúp em đỡ ngu.
\end{enumerate}

\section{25 câu về lười mà hay chống chế}
\begin{enumerate}[leftmargin=*,resume]
\item Lười mà biết mình lười còn cứu được, lười mà cứ gọi là chờ cảm hứng thì khó.
\item Trì hoãn là cách rất lịch sự để tự phá mình.
\item Đừng gọi là chưa sẵn sàng khi thật ra em chỉ đang sợ bắt đầu.
\item Cái điện thoại không giết tương lai em, nhưng việc em không thắng nổi nó thì có.
\item Ai cũng mệt, khác nhau ở chỗ có người vẫn làm phần việc phải làm.
\item Đừng nhận mình áp lực nếu em còn tiêu thời gian như người rảnh lắm.
\item Nhiều đứa than thiếu thời gian nhưng lại rất dư thời gian cho mấy thứ vô bổ.
\item Tự hứa rồi tự nuốt lời nhiều quá thì đừng hỏi sao em coi thường chính mình.
\item Cảm hứng là món quà, kỷ luật mới là thứ nuôi em lớn.
\item Bản thân yếu mà cứ mong hoàn cảnh thuận là đang nằm mơ ban ngày.
\item Một ngày lười không giết ai, nhưng lười có lý luận thì giết dần tương lai.
\item Đừng nói mình bận khi việc quan trọng nhất toàn bị em đẩy xuống cuối.
\item Người mệt thật thường tìm cách tối giản, người lười thường tìm cách ngụy biện.
\item Thích thành quả nhưng ghét quá trình là kiểu trẻ con rất dai.
\item Muốn giỏi nhanh thường là dấu hiệu của cái tôi đói hơn năng lực.
\item Thức khuya không chứng minh em chăm, có khi chỉ chứng minh em quản mình kém.
\item Đừng khoe bận tới mức không ngủ được nếu phần lớn là em tự phá lịch mình.
\item Cứ đợi mai rồi mai nữa thì đời cũng đợi em tụt lại thôi.
\item Lười đâu đáng sợ bằng lười nhưng nói chuyện như mình đang hy sinh lớn lắm.
\item Nỗ lực giả là làm nhiều việc nhỏ để khỏi động vào việc khó.
\item Càng trì hoãn càng dễ ghét chính cái việc mình chưa làm.
\item Đừng nói số em xui khi em còn chưa kỷ luật nổi giờ ngủ.
\item Không ai cứu nổi người chuyên tự thương hại để hợp thức hóa sự lười.
\item Có những giấc mơ chết không vì nghèo mà vì chủ nhân của nó quen dễ dãi.
\item Đã yếu mà còn thích tự tha là công thức tụt dài hạn.
\end{enumerate}

\section{25 câu về tỉnh ngộ muộn vẫn còn hơn không}
\begin{enumerate}[leftmargin=*,resume]
\item Nhận ra mình dở hơi muộn vẫn đỡ hơn sống dở mà tưởng ngầu.
\item Có những ngày xấu hổ chính là ngày đầu em bắt đầu lớn.
\item Tỉnh ngộ không đẹp, thường nó đau và hơi quê.
\item Cái tát của sự thật thường đỡ hại hơn cái vuốt ve của ảo tưởng.
\item Đừng tiếc hình ảnh cũ của mình nếu nó đang kéo em lùi.
\item Người trưởng thành không phải người chưa từng ngu, mà là người dám thôi ngu.
\item Nhiều cú ngã đáng tiền vì nó làm em bớt ảo.
\item Bị chạm tự ái đôi khi là học phí của việc mở mắt.
\item Đừng sợ thừa nhận mình sai, chỉ sợ sai mà cố giữ vai đúng.
\item Sống thật bắt đầu từ lúc em thôi diễn cho chính mình xem.
\item Không ai khá lên bằng cách bảo vệ mãi phiên bản tệ của bản thân.
\item Cái khó nhất không phải biết mình sai, mà là bỏ cái lợi nhỏ do cái sai đem lại.
\item Muốn đổi đời đôi khi phải đổi cái kiểu nghĩ đã nuôi em lớn lên.
\item Sự thật không cần dễ nghe mới đúng.
\item Tỉnh táo luôn đắt hơn tự ru ngủ.
\item Có những người ghét em sau khi em thay đổi, vì họ quen hưởng lợi từ bản cũ của em.
\item Càng tự ái càng học chậm.
\item Đừng mong cuộc nói chuyện nào cũng nhẹ nhàng nếu điều cần nghe vốn không nhẹ.
\item Có những lúc bị chê không phải bị hạ thấp, mà đang được kéo về đúng mặt đất.
\item Thấy mình còn nhiều chỗ dở là dấu hiệu tốt, nghĩa là mắt em bớt mù.
\item Bỏ cái tôi xuống một chút, đầu óc sẽ có chỗ để lớn lên.
\item Nhận lỗi không làm em thấp đi, chỉ làm em bớt giả.
\item Người dám sửa mình thường bớt nói lớn về bản thân.
\item Không ai thoát khỏi giai đoạn ngu đầu đời, vấn đề là em ở đó bao lâu.
\item Muộn còn hơn là chết dí trong cái sai mà vẫn cười tỉnh.
\end{enumerate}

\section{25 câu về giá trị thật của học hành}
\begin{enumerate}[leftmargin=*,resume]
\item Học để có nghề chỉ là một nửa, nửa còn lại là học để có não mà sống.
\item Người có chữ chưa chắc tử tế, nhưng người cố tình anti tri thức thì khó sâu.
\item Học giúp em kiếm tiền, nhưng quan trọng hơn là giúp em đỡ bán rẻ mình.
\item Càng hiểu đời càng bớt thích nói to mấy câu đơn giản hóa mọi thứ.
\item Học là cách âm thầm xây giá trị lúc chưa ai nhìn thấy.
\item Kiến thức không luôn cứu em khỏi sai, nhưng nó giúp em sai đỡ ngu.
\item Có chữ không làm em sang chảnh, nó làm em đỡ thô với đời.
\item Người biết nghĩ thường ít bị dẫn bằng mấy câu rất dễ kích động.
\item Học không chỉ để trả lời, còn để biết câu hỏi nào đáng trả lời.
\item Một người trẻ mà ghét nghĩ là tự đóng cửa tương lai của mình.
\item Càng nhiều công cụ thông minh, càng không được phép sống ngu yên tâm.
\item Học không làm em thành thiên tài, nhưng giúp em bớt thành nạn nhân.
\item Không phải cái gì có thể tra mạng cũng thay được năng lực hiểu thật.
\item Học là việc âm thầm, còn hậu quả của thiếu học thường rất ồn.
\item Người chịu học thường bớt cay cú với thành công của người khác.
\item Có kiến thức rồi em mới biết mình cần học tiếp chỗ nào.
\item Đừng coi thường nền tảng, đa số cú sập đều bắt đầu từ phần móng.
\item Học để thấy thế giới lớn hơn cái khuôn em đang sống.
\item Càng đọc nhiều càng bớt tin mấy câu rất chắc mà rất ngu.
\item Tri thức không làm em hết cô đơn, nhưng cho em ngôn ngữ để hiểu nỗi cô đơn đó.
\item Ai cũng muốn sống có giá trị, nhưng không phải ai cũng chịu học để tạo giá trị.
\item Học thật chưa chắc được khen sớm, nhưng thường được đời trả sau.
\item Muốn bớt phụ thuộc thì phải tăng năng lực.
\item Càng ít nội lực càng dễ phải sống nhờ hào quang vay mượn.
\item Học nghiêm túc là cách tôn trọng phiên bản tương lai của mình.
\end{enumerate}

\chapter{100 câu cục súc về lười, trì hoãn và tự phá mình}
\markboth{100 câu cục súc về lười, trì hoãn và tự phá mình}{100 câu cục súc về lười, trì hoãn và tự phá mình}

\section{25 câu bóc trần kiểu lười có văn}
\begin{enumerate}[leftmargin=*]
\item Đừng gọi mình cầu toàn nếu em toàn trì hoãn tới lúc làm đại.
\item Nói đang tích năng lượng nghe hay, nhưng nhiều khi chỉ là lười đẹp lời.
\item Mệt thật cần nghỉ, lười thường cần cớ.
\item Không thiếu thời gian, em đang thiếu quyết tâm cho cái thật sự quan trọng.
\item Nhiều người rất siêng chạy trốn việc khó.
\item Nói chưa đúng thời điểm là cách trì hoãn được đóng gói sang hơn.
\item Tự thưởng quá nhiều trước khi làm xong là thói quen rất phá đời.
\item Cứ nói áp lực trong khi nửa ngày chìm trong điện thoại thì nghe không thuyết phục lắm.
\item Lười kéo dài sẽ đẻ ra tự ti, rồi tự ti lại làm em lười thêm.
\item Trốn việc một chút thấy nhẹ, trốn lâu thành người yếu.
\item Càng né càng sợ, càng sợ càng né, rồi em tự gọi đó là tính cách.
\item Có những người không thất bại vì dở, mà vì đợi quá lâu mới bắt đầu.
\item Đừng than thiếu động lực khi em chưa từng dựng nổi nề nếp.
\item Kế hoạch đẹp mà tay không nhúc nhích thì cũng chỉ là tranh treo tường.
\item Nhiều đứa mê xem video phát triển bản thân hơn là phát triển bản thân thật.
\item Cái cảm giác mình sắp làm lớn thường rất dễ gây nghiện với người chưa làm gì cả.
\item Lười không luôn ồn, đôi khi nó mặc áo bận rộn.
\item Bận giả là làm mười chuyện lặt vặt để khỏi làm một chuyện quan trọng.
\item Không phải em thiếu khả năng, nhiều lúc em chỉ quá quen dễ dãi.
\item Muốn đổi mà chưa muốn khó thì thường chỉ đổi status.
\item Thói quen bé nhưng lặp lại mỗi ngày mạnh hơn quyết tâm bốc đồng theo tuần.
\item Tự hứa ít thôi, làm nhiều hơn đi, đời không ăn lời phát biểu.
\item Cứ nói mai bắt đầu là đang cho bản thân uống thuốc mê.
\item Trì hoãn giỏi thì đến lúc chữa hậu quả cũng phải giỏi.
\item Người tự quản kém rất thích đổ cho hoàn cảnh thiếu cảm hứng.
\end{enumerate}

\section{25 câu về cái giá của việc chần chừ}
\begin{enumerate}[leftmargin=*,resume]
\item Cái giá đầu tiên của trì hoãn là mất tự trọng.
\item Càng để việc treo càng nặng đầu, nhưng em vẫn thích tự treo thêm.
\item Làm muộn một lần là sự cố, làm muộn thành nếp là tính cách.
\item Nhiều giấc ngủ không yên đến từ những việc em biết mình đang né.
\item Chần chừ lâu ngày biến chuyện nhỏ thành nỗi sợ lớn.
\item Cứ để việc dồn là đang vay tương lai với lãi rất cao.
\item Muộn giờ nhiều rồi em sẽ thấy mình tin vào bản thân ngày càng ít.
\item Việc chưa làm xong rất biết cách hút hết năng lượng nền của em.
\item Nỗi mệt vì dồn việc thường do chính em sản xuất.
\item Trì hoãn không chỉ cướp thời gian, nó cướp cả sự yên ổn.
\item Làm gấp liên tục sẽ khiến em lầm tưởng hỗn loạn là trạng thái bình thường.
\item Cứ để áp lực đuổi mới chạy thì cả đời chỉ quen chạy trốn.
\item Không ai xây được đời bền bằng những pha cứu phút cuối kéo dài nhiều năm.
\item Khó chịu nhất không phải việc khó, mà là biết mình hoàn toàn có thể làm sớm hơn.
\item Tự kéo mình vào góc rồi than số phận ép là kiểu rất thiếu công bằng với đời.
\item Mỗi lần em bỏ qua một cảnh báo nhỏ là em đang nuôi một vấn đề lớn hơn.
\item Chần chừ làm khả năng tập trung của em nát dần mà ít ai để ý.
\item Nợ việc còn làm em cáu bẳn với cả người không có lỗi.
\item Trì hoãn lâu sẽ dạy não em rằng né tránh luôn có thể chấp nhận.
\item Một cuộc đời cứ sống bằng hạn chót rất khó có chiều sâu.
\item Những việc bị em né hôm nay thường quay lại với bộ mặt xấu hơn ngày mai.
\item Càng quen sống sát giờ càng khó nghĩ dài.
\item Nhiều người mất cơ hội không vì kém, mà vì mọi thứ trong tay họ luôn đến trễ.
\item Nếu cái gì cũng chờ nước tới chân mới nhảy, sớm muộn có bữa không còn chân để nhảy.
\item Chần chừ là cách âm thầm nói với bản thân rằng việc mình làm không đáng được ưu tiên.
\end{enumerate}

\section{25 câu về kỷ luật thô nhưng hiệu quả}
\begin{enumerate}[leftmargin=*,resume]
\item Muốn bớt trì hoãn thì đừng mặc cả với việc phải làm.
\item Kỷ luật bắt đầu từ mấy chuyện rất nhỏ mà em hay xem thường.
\item Đừng đợi có hứng mới ngồi vào bàn, ngồi vào bàn rồi hứng mới biết đường tới.
\item Làm trước phần khó một chút, em sẽ bớt sợ cả ngày.
\item Tự quản được sáng sớm thì nhiều chuyện trong đời tự dễ hơn.
\item Lịch không cứu em, nhưng lịch rõ sẽ bớt cho em diễn với chính mình.
\item Việc lớn luôn chết vì những lần nhượng bộ nhỏ lặp lại.
\item Bỏ điện thoại xa tay là bước văn minh nhất của nhiều người trẻ.
\item Càng dễ bị xao nhãng càng phải thiết kế môi trường tử tế hơn.
\item Đầu óc không sinh ra để liên tục chống lại cám dỗ miễn phí em tự để cạnh mình.
\item Kỷ luật không vui, nhưng hậu quả của vô kỷ luật còn vui hơn chắc.
\item Người mạnh không phải người không muốn nghỉ, mà là người biết nghỉ xong rồi làm tiếp.
\item Đừng chờ một ngày hoàn hảo, hãy cứu lấy một giờ còn dùng được.
\item Mỗi lần em làm đúng dù không thích là em đang xây cơ bắp tinh thần.
\item Có nề nếp rồi mới biết năng lượng mình từng bị thất thoát vô lý cỡ nào.
\item Làm đều nghe chán, nhưng chính cái chán đó nuôi thành quả.
\item Kỷ luật là món đắng ban đầu nhưng ít tác dụng phụ về sau.
\item Muốn tự do lâu dài thì chấp nhận kỷ luật ngắn hạn.
\item Người trẻ mạnh không phải người nói hay, mà là người giữ được lời với bản thân.
\item Kỷ luật là cách bảo vệ mơ ước khỏi tâm trạng thất thường.
\item Không ai quản em vẫn làm là lúc bản lĩnh bắt đầu hình thành.
\item Thói quen đúng nhìn nhỏ, nhưng nó gánh cả tương lai rất to.
\item Muốn khác số đông thì đừng sống theo kiểu tiện đâu làm đó.
\item Kỷ luật không làm em đặc biệt liền, nó làm em đáng tin trước đã.
\item Đáng tin với chính mình là nền của mọi kiểu trưởng thành khác.
\end{enumerate}

\section{25 câu về thôi tự phá mình}
\begin{enumerate}[leftmargin=*,resume]
\item Đừng đợi ai cứu nếu em vẫn là người phá mình nhanh nhất.
\item Có những người không cần kẻ thù vì chính nếp sống của họ đã đủ ác.
\item Tự coi thường mình lâu ngày là kiểu tự hại rất âm.
\item Ăn ngủ thất thường rồi than tâm trạng bất ổn nghe cũng hợp lý theo cách buồn cười.
\item Cái miệng nói thương bản thân nhưng lịch sống đang hành hạ bản thân.
\item Đừng đòi đầu óc sáng nếu cơ thể em bị đối xử như đồ dùng rẻ tiền.
\item Tự phá mình thường núp dưới vỏ bọc vui là chính mình.
\item Không phải cái gì em thích cũng đang tốt cho em.
\item Yêu bản thân không phải để chiều mình vô điều kiện.
\item Nhiều người gọi đó là tận hưởng, thật ra là tiêu hao không kiểm soát.
\item Tâm trạng xấu không phải lúc nào cũng tại đời, có khi tại nếp sống em tự thả.
\item Tự do không có nghĩa là ngày nào cũng làm điều hại mình.
\item Nếu em luôn chọn thứ dễ nhất, đừng ngạc nhiên khi đời đối xử với em rất cứng.
\item Người khôn không phải người chưa từng tự phá, mà là người biết dừng đúng lúc.
\item Đừng biến sự buông thả thành cá tính.
\item Có những thói quen em gọi là bình thường, nhưng tương lai của em sẽ gọi nó là nguyên nhân.
\item Muốn bớt tệ thì trước hết phải thôi bênh cái tệ của mình.
\item Tự chăm mình là chuyện thật, tự nuông hỏng mình là chuyện khác.
\item Thấy mình đang trôi mà còn khoái cảm giác trôi là rất nguy hiểm.
\item Người tỉnh táo sẽ sớm nhận ra vài thú vui đang âm thầm rút máu mình.
\item Đừng để vài phút sướng cướp vài năm đáng sống.
\item Tự phá dễ vì hậu quả không đến ngay.
\item Cái gì làm em vui tức thời nhưng rỗng dần về sau thì nên dè chừng.
\item Không ai lớn lên tử tế bằng cách suốt ngày làm điều khiến mình xấu đi.
\item Thương mình thật là bớt tiếp tay cho những thứ đang làm mình nát.
\end{enumerate}

\chapter{100 câu cục súc về điểm số, sĩ diện và áp lực}
\markboth{100 câu cục súc về điểm số, sĩ diện và áp lực}{100 câu cục súc về điểm số, sĩ diện và áp lực}

\section{25 câu về điểm thấp và sự thật khó nuốt}
\begin{enumerate}[leftmargin=*]
\item Điểm thấp không làm em vô dụng, nhưng nó đang báo chỗ em hổng.
\item Cái đau nhất của điểm thấp là nó không chiều ảo tưởng của em.
\item Đề khó có thể là thật, nhưng phần em chưa biết cũng là thật.
\item Cứ chửi đề mãi không làm bài em đúng thêm được câu nào.
\item Có những cú điểm thấp chỉ để dạy em bớt tưởng mình ổn.
\item Đừng bảo học tài thi phận khi em còn chưa học tới nơi.
\item Xem điểm mà tim thắt lại là vì em biết mình đã bỏ qua gì.
\item Điểm số không nói hết giá trị của em, nhưng nó nói được phần trách nhiệm học tập gần đây.
\item Đừng lấy cảm giác hiểu bài để thế chỗ cho năng lực làm bài.
\item Điểm kém không đáng sợ bằng cái kiểu né xem bài để khỏi đau.
\item Sai một bài chưa nát, trốn nhìn cái sai mới dễ nát.
\item Điểm thấp là thông tin, đừng biến nó thành bản án rồi nằm im.
\item Cảm giác bị hụt khi xem điểm là tín hiệu để sửa, không phải để tự chôn.
\item Có người thích mở điểm thật nhanh, có người sợ mở điểm, ai cũng đang đối diện bản thân theo cách khác nhau.
\item Nếu em chỉ muốn điểm đẹp mà không chịu nhìn lỗi xấu thì hơi vô lý.
\item Điểm số không công bằng tuyệt đối, nhưng nó vẫn công bằng hơn nhiều lời tự khen của bản thân.
\item Đừng coi một bài thấp là bằng chứng em hết cứu.
\item Người mạnh không phải người luôn điểm cao, mà là người dám soi bài thấp rất kỹ.
\item Điểm thấp không giết tương lai, thói quen né trách nhiệm mới làm điều đó.
\item Càng giấu điểm càng giấu luôn cơ hội sửa.
\item Đừng ghét con số, hãy ghét sự hổng mà con số đang chỉ ra.
\item Có những bài điểm thấp rất đáng quý vì nó kéo em về thực lực thật.
\item Chưa giỏi thì thừa nhận, đỡ mệt hơn cố giữ mặt.
\item Nỗi đau của điểm thấp ngắn hơn hậu quả của việc sống ảo về năng lực.
\item Một bài tệ không định nghĩa em, nhưng thái độ sau bài tệ thì có.
\end{enumerate}

\section{25 câu về sĩ diện học đường}
\begin{enumerate}[leftmargin=*,resume]
\item Nhiều đứa không sợ dốt, chỉ sợ người khác thấy mình dốt.
\item Sĩ diện là thứ làm em im miệng đúng lúc cần hỏi.
\item Muốn khá lên mà không muốn quê là bài toán rất trẻ con.
\item Đừng cố làm người hiểu hết nếu em còn đang mù nửa bài.
\item Sĩ diện giữ được mặt vài phút, thường phá tiến bộ vài tháng.
\item Lên bảng đứng hình không đáng nhục bằng việc ở dưới cứ giả bộ hiểu.
\item Cái tôi lớn quá thì chỗ chứa kiến thức sẽ ít đi.
\item Nhiều người học chậm vì quá bận diễn vai ổn.
\item Đừng ngại hỏi cái cơ bản, ngoài đời người ta chỉ ngại người cứng đầu mà rỗng.
\item Có khi em không cần thêm thông minh, em cần bớt sĩ diện.
\item Giữ hình tượng quá lâu thì chẳng còn sức giữ sự thật.
\item Trưởng thành là lúc dám nói mình chưa biết mà không tự thấy thấp.
\item Sĩ diện làm người ta thích cái vỏ đúng hơn cái lõi tiến bộ.
\item Cứ muốn giỏi sẵn trước mặt người khác thì không bao giờ chấp nhận đủ thời gian để học.
\item Mất mặt một chút đôi khi là giá rẻ để mở đầu cho năng lực thật.
\item Có những người nhìn rất tự tin, thật ra chỉ đang phòng thủ bằng thái độ.
\item Đừng cười người hỏi ngu, kẻ đáng sợ hơn là người ngu mà không hỏi.
\item Sĩ diện làm em dị ứng với phản hồi, rồi tự khóa đường sửa.
\item Tự trọng khác sĩ diện, một thứ giúp em lớn, một thứ giữ em nhỏ mãi.
\item Càng dễ quê càng nên tập đứng ở chỗ chưa biết.
\item Muốn người khác tôn trọng thì trước hết phải tôn trọng sự thật về trình độ của mình.
\item Mặt mũi không nuôi nổi tương lai, năng lực mới nuôi được.
\item Người dám nhận dở thường tiến nhanh hơn người cứ diễn giỏi.
\item Bị lộ cái chưa biết không đáng sợ bằng cả đời sống trong cái chưa biết đó.
\item Cái tôi đói khen rất hay ghét câu nói thật.
\end{enumerate}

\section{25 câu về so sánh và áp lực}
\begin{enumerate}[leftmargin=*,resume]
\item So nỗi khổ hậu trường của mình với highlight của người khác là tự hành xác.
\item Đừng nhìn người khác đang bay mà quên coi chân mình đang đứng ở đâu.
\item Áp lực không xấu hết, nhưng sống chỉ để vượt người khác thì rất mệt.
\item So sánh quá liều sẽ biến cả thành công của người khác thành nhát dao vào lòng mình.
\item Không phải ai hơn em cũng là kẻ đe dọa em.
\item Có người hơn em vì họ bền hơn, không phải vì họ được trời thương hơn.
\item Ghen tị không làm em tiến, cùng lắm chỉ làm em thức khuya hơn.
\item Thấy người khác giỏi mà cay là dấu hiệu em chưa ổn với chính mình.
\item So sánh có ích khi nó chỉ ra việc cần làm, vô ích khi nó chỉ châm thêm tự ti.
\item Áp lực lớn nhất thường đến từ câu em tự nói với mình sau khi thấy ai đó giỏi hơn.
\item Đừng để bảng điểm biến bạn bè thành đối thủ tưởng tượng suốt ngày.
\item Người nào cũng có đoạn lên nhanh, đoạn đứng lại, đoạn tụt; em chỉ hay thấy đoạn họ lên.
\item So sánh thiếu dữ liệu là trò rất tàn nhẫn với bản thân.
\item Càng nhìn ngang nhiều quá càng khó nhìn sâu vào việc mình phải làm.
\item Áp lực từ gia đình đau, nhưng áp lực tự dựng trong đầu đôi khi còn đau hơn.
\item Không ai bắt em thành bản sao điểm cao của người khác.
\item Cái đích của em phải thật, đừng là ảnh chụp thành công của thiên hạ.
\item Cứ nhìn bạn học hơn mình mà quên hôm qua mình là cách sống rất mỏi.
\item Người khôn dùng thành công của người khác làm dữ liệu, không làm dao cứa mình.
\item Áp lực có thể kéo em đi, nhưng đừng để nó dắt em như trâu.
\item Mỗi người một quỹ thời gian, nền tảng và nỗi lo, so kiểu bằng phẳng là sai từ gốc.
\item Nếu em chỉ vui khi hơn người, đời em sẽ ít ngày vui lắm.
\item Đừng biến việc học thành cuộc chiến giữ mặt với cả thế giới.
\item Cần phấn đấu, không cần sống như đang bị truy đuổi bởi mọi bảng xếp hạng.
\item Cái cần vượt lâu dài là bản thân lỏng tay của em, không phải cái bóng của người khác.
\end{enumerate}

\section{25 câu về đối diện kết quả cho đàng hoàng}
\begin{enumerate}[leftmargin=*,resume]
\item Xem kết quả đi rồi mới có quyền buồn đúng chỗ.
\item Trốn điểm không làm điểm đổi, chỉ làm em yếu thêm.
\item Càng né kết quả càng ít cơ hội sửa sai còn nóng.
\item Người có bản lĩnh không thích sự thật hơn người khác, họ chỉ chịu nghe nó hơn.
\item Đối diện kết quả là bước đầu của mọi lần làm lại cho ra hồn.
\item Đừng chỉ đem điểm về nhà, hãy đem cả kế hoạch sửa.
\item Học cách chịu trách nhiệm với con số của mình trước khi đòi đời tôn trọng ước mơ của mình.
\item Có thể buồn, có thể tức, nhưng cuối cùng vẫn phải ngồi xuống sửa bài.
\item Kết quả xấu không đòi em gục, nó đòi em tỉnh.
\item Người ta nể em không phải vì em không từng thấp, mà vì em thấp rồi vẫn ngoi lên được.
\item Đừng biến một bài tệ thành câu chuyện bi kịch dài tập.
\item Càng nói mình xui nhiều quá càng khó nhìn ra thứ cần thay đổi.
\item Nhận trách nhiệm không phải tự đánh mình, mà là giành lại quyền sửa.
\item Nếu em chưa dám nhìn kết quả của mình thì ước mơ lớn nghe hơi non.
\item Mạnh không phải là không đau, mà là đau xong vẫn sửa tiếp.
\item Một người trẻ tử tế nên học cách thừa nhận kết quả của chính tay mình làm ra.
\item Đừng đợi ai an ủi đủ mới chịu bắt đầu lại.
\item Ngồi khóc lâu trước bài điểm thấp không giúp em đúng thêm dòng nào.
\item Càng rõ nguyên nhân càng bớt bất lực.
\item Học sinh lớn lên nhanh khi biết biến một bài tệ thành bản đồ cho lần sau.
\item Đừng so nỗi đau của em với ai, hãy so hành động sửa của em với hôm qua.
\item Kết quả không đẹp là lúc nhân cách thi thử luôn.
\item Dám báo, dám nhận, dám sửa là ba bước ít người trẻ làm gọn được.
\item Có trách nhiệm rồi, điểm số dù chưa đẹp vẫn còn hy vọng.
\item Đời không nể tiếng than, đời nể sức bật sau cú thấp.
\end{enumerate}

\chapter{100 câu cục súc về tình bạn, lòng người và cô độc}
\markboth{100 câu cục súc về tình bạn, lòng người và cô độc}{100 câu cục súc về tình bạn, lòng người và cô độc}

\section{25 câu về bạn bè và sự chọn lọc}
\begin{enumerate}[leftmargin=*]
\item Bạn tốt không làm em thấy vui mọi lúc, nhưng thường làm em bớt ngu.
\item Đừng gọi tất cả là bạn, nhiều người chỉ đi cùng đoạn tiện.
\item Chơi với ai lâu ngày sẽ mượn luôn cách nghĩ của người đó.
\item Bạn bè kéo em xuống mà em còn gọi là vui thì hơi dở.
\item Không phải ai hiểu nỗi buồn của em cũng muốn em tốt lên.
\item Có đứa ở cạnh em vì em dễ lợi dụng, không phải vì quý.
\item Tình bạn bền hiếm khi cần diễn nhiều.
\item Đừng giữ một nhóm chỉ vì sợ cô đơn.
\item Có người cười với em rất tươi nhưng không hề muốn em đi xa hơn họ.
\item Bạn thật không thích em yếu, họ chỉ không bỏ em lúc em yếu.
\item Hợp nói chuyện không đồng nghĩa hợp đi đường dài.
\item Người luôn khiến em thấy mình tệ đi không đáng có vé ngồi lâu trong đời em.
\item Đừng nhầm sự quen mặt với sự thân thiết.
\item Nhiều mối quan hệ chết không vì cãi nhau, mà vì lớn lên khác hướng.
\item Chơi với người hay chê mọi cố gắng của em lâu ngày em cũng mất lửa.
\item Bạn bè tốt không thay đời em, nhưng thay quỹ đạo nghĩ của em rất mạnh.
\item Đừng giao phần mong manh nhất của mình cho người chỉ biết đùa.
\item Không phải nhóm đông là đỡ lạc lõng.
\item Quý ai đó không có nghĩa là phải nuốt hết mọi kiểu đối xử của họ.
\item Bạn bè đàng hoàng sẽ tôn trọng ranh giới của em.
\item Có những cuộc chơi rút máu tinh thần hơn cả một bài thi khó.
\item Người nào chỉ nhớ tới em lúc cần thường nên để ở đúng vị trí cần thôi.
\item Bạn bè cũng cần chọn như chọn đường, đi sai lâu mệt lắm.
\item Ở cạnh người tử tế lâu ngày em sẽ bớt thích mấy trò nông cạn.
\item Chọn bạn là đang chọn phần nào con người tương lai của mình.
\end{enumerate}

\section{25 câu về thất vọng và lòng người}
\begin{enumerate}[leftmargin=*,resume]
\item Đừng đau quá vì người ta thay đổi, nhiều khi họ chỉ lộ ra thôi.
\item Lòng người không xấu hết, nhưng ngây thơ quá thì em tự rước đau.
\item Tin người là cần, tin mù là tự nộp mình.
\item Có người tử tế cho tới khi quyền lợi của họ bị đụng.
\item Thất vọng lớn thường đến từ kỳ vọng em tự thêu quá đẹp.
\item Không ai có nghĩa vụ giữ nguyên vai tốt trong câu chuyện em tự viết về họ.
\item Bị phản bội đau vì nó đập vào chỗ mình đã từng tin rất thật.
\item Đừng vì một vài người tệ mà hóa đá với tất cả, nhưng cũng đừng vì cần người mà mở cửa bừa.
\item Lòng người đổi chiều nhanh hơn nhiều lời hứa.
\item Có những người chỉ tốt khi em chưa hơn họ.
\item Bài học về lòng người thường đắt vì học bằng tin tưởng thật.
\item Bớt kể hết về mình cho người mới quen là một kiểu khôn cần thiết.
\item Người nào luôn cần em nhưng chưa từng thật sự hiện diện cho em thì nên xem lại.
\item Có người nói thương em, nhưng họ chỉ thương cảm giác được em cần.
\item Đừng sốc khi ai đó ích kỷ, hãy sốc nếu sau nhiều dấu hiệu em vẫn giả vờ không thấy.
\item Tin quá nhanh thường là do mình đói sự đồng cảm.
\item Nỗi đau lớn nhất không phải người ta xấu, mà là mình từng tự tin họ không như vậy.
\item Càng lớn càng hiểu không phải ai cười với mình cũng cùng phe.
\item Khôn lên không phải mất niềm tin, mà là đặt niềm tin đúng chỗ.
\item Có người bước vào đời em để dạy một bài học chứ không để ở lại.
\item Đừng trách tim mềm nếu sau cùng đầu em biết rút kinh nghiệm.
\item Người chân thành thường không cần phô diễn quá nhiều về sự chân thành của họ.
\item Học nhìn người không để nghi ngờ cả thế giới, mà để đỡ tự nộp mình thêm lần nữa.
\item Sau vài cú vỡ mộng, em sẽ biết bình tĩnh quan sát quý hơn cảm giác được yêu quý ngay.
\item Lòng người là bài thi không ai chấm trước đáp án cho em.
\end{enumerate}

\section{25 câu về cô đơn và tự đứng bằng chân mình}
\begin{enumerate}[leftmargin=*,resume]
\item Cô đơn không luôn vì thiếu người, nhiều khi vì thiếu nơi được hiểu.
\item Ở giữa đám đông mà lạc vẫn là một kiểu cô đơn rất thật.
\item Không phải ai cũng có phước được gặp người nói đúng thứ mình cần nghe.
\item Cô đơn không giết em nhanh, nhưng nếu không hiểu nó, nó mài em rất đều.
\item Cần người là bình thường, lệ thuộc vào việc luôn có người mới mệt.
\item Có những đêm chỉ cần tự chịu được mình đã là một loại trưởng thành.
\item Học ở một mình không vui, nhưng cần.
\item Người càng không dám một mình càng dễ bấu vào sai người.
\item Cô đơn nhiều lúc chỉ là khoảng trống giữa con người thật và vai diễn ngoài mặt.
\item Đừng vội lấp nỗi trống bằng bất kỳ ai đến đúng lúc em yếu.
\item Có những giai đoạn không ai hiểu em, nhiệm vụ của em vẫn là đừng bỏ em.
\item Biết tự dỗ mình là kỹ năng không sexy nhưng rất cứu mạng.
\item Một người không chịu nổi im lặng thường đang tránh một điều gì đó trong mình.
\item Cô đơn cho em thấy em đang đối thoại với bản thân kiểu nào.
\item Đừng biến nỗi cô đơn thành lý do để chọn đại một mối quan hệ.
\item Ở một mình mà bình yên tốt hơn ở cạnh người làm em cạn năng lượng.
\item Trưởng thành là lúc em ngừng dùng người khác để tạm quên mình.
\item Cần được yêu là thật, nhưng trước đó cần bớt tự bỏ rơi chính mình.
\item Nhiều người không hề cô đơn, họ chỉ đang cai cảm giác được chú ý liên tục.
\item Một trái tim rối rất dễ tưởng tiếng ồn là sự chữa lành.
\item Cô đơn đúng cách có thể dạy em nhiều điều mà hội nhóm ồn ào không dạy được.
\item Đừng sợ khoảng lặng, nhiều câu trả lời chỉ hiện ra trong đó.
\item Có những ngày chỉ cần không phản bội bản thân là đủ tốt.
\item Tự đứng được rồi em sẽ bớt sợ mất những người vốn không thuộc về mình.
\item Một mình không dễ chịu, nhưng nó làm em rõ mình hơn nhiều cuộc vui giả.
\end{enumerate}

\section{25 câu về giữ mình giữa lòng người}
\begin{enumerate}[leftmargin=*,resume]
\item Giữ lòng tốt nhưng đừng bỏ đầu óc.
\item Mềm lòng là đẹp, mềm nguyên tắc là tự hại.
\item Đừng để vài lần thất vọng làm em hóa thành người cay độc.
\item Sống tử tế không có nghĩa là đưa ai cũng quyền làm đau mình.
\item Đàng hoàng là biết cho đi, nhưng cũng biết dừng đúng lúc.
\item Có ranh giới không làm em lạnh, nó làm em lành.
\item Đừng cố giải thích giá trị của mình với người chỉ muốn lợi dụng.
\item Bớt mong được mọi người thích thì em sẽ dễ sống đúng hơn.
\item Không cần thắng trong mọi mối quan hệ, chỉ cần không thua chính mình.
\item Tha thứ là lựa chọn của em, quay lại như cũ hay không là chuyện khác.
\item Có người chỉ hiểu sự tử tế của em như một cơ hội để leo lên đầu.
\item Đừng biến lòng thương thành chiếc khóa trói chân mình.
\item Bảo vệ bình yên không ích kỷ, nó là trách nhiệm với đời sống tinh thần của em.
\item Càng biết mình quý ở đâu càng bớt cần giữ những người coi thường mình.
\item Có những cánh cửa đóng lại không phải mất mát, mà là lọc rác.
\item Đừng cố làm nơi trú cho người cứ thích đốt nhà người khác.
\item Người tử tế thật thường cũng rất rõ ràng khi cần.
\item Không phải ai khổ cũng có quyền làm khổ em.
\item Đừng đem tim mình làm trạm cứu hộ cho tất cả.
\item Sống ấm nhưng không mềm nhũn là một bản lĩnh lớn.
\item Chọn bình yên không phải hèn, nhiều khi là khôn.
\item Bớt chứng minh mình tốt cho người không cần điều tốt.
\item Em không cần cứu mọi người để chứng minh mình đáng quý.
\item Giữ mình trước lòng người không làm em mất nhân tính, nó giữ nhân tính của em khỏi bị bào mòn.
\item Tử tế có xương sống mới đi đường dài được.
\end{enumerate}

\chapter{100 câu cục súc về gia đình, kỳ vọng và tổn thương}
\markboth{100 câu cục súc về gia đình, kỳ vọng và tổn thương}{100 câu cục súc về gia đình, kỳ vọng và tổn thương}

\section{25 câu về kỳ vọng của gia đình}
\begin{enumerate}[leftmargin=*]
\item Gia đình thương em không có nghĩa là gia đình luôn hiểu đúng em.
\item Nhiều áp lực mang tên lo cho con nhưng cách đưa lại như đang siết cổ con.
\item Kỳ vọng không xấu, nhưng kỳ vọng mà thiếu lắng nghe thì dễ thành gánh.
\item Có nhà nói yêu bằng hy sinh, có nhà nói yêu bằng kiểm soát, không phải lúc nào hai thứ đó cũng giống nhau.
\item Người lớn đôi khi bắt em sống giấc mơ dang dở của họ rồi gọi đó là định hướng.
\item Thương con mà chỉ biết so sánh là cách làm con mòn rất nhanh.
\item Không phải cha mẹ nào cũng biết nói điều đúng theo cách lành.
\item Con cái im không phải luôn vì ngoan, có khi vì mệt nói cũng không ai nghe.
\item Gia đình có thể là chỗ dựa, cũng có thể là nơi phát sinh nỗi sợ lớn nhất của một người trẻ.
\item Nhiều vết xước trong lòng trẻ con mang giọng người lớn rất quen.
\item Đừng vì biết ơn mà coi nhẹ nỗi đau của mình.
\item Hiếu thảo không đồng nghĩa với phải nuốt hết mọi tổn thương rồi cười.
\item Có những đứa trẻ lớn lên trong nhà đủ ăn nhưng thiếu chỗ được hiểu.
\item Lời chê từ người lạ đau một lúc, từ gia đình đau lâu hơn.
\item Người trẻ nào cũng muốn được nhìn như một con người, không chỉ như bảng thành tích di động.
\item Gia đình áp lực thường tạo ra hai kiểu con: hoặc rất ngoan ngoài mặt, hoặc rất lì bên trong.
\item Không ít người giỏi bên ngoài nhưng kiệt sức vì mãi cố đáng giá trong mắt nhà mình.
\item Kỳ vọng quá cao mà không có kỹ năng đồng hành chỉ sinh ra khoảng cách.
\item Yêu thương thiếu tinh tế đôi khi để lại hậu quả rất dài.
\item Có những đứa trẻ không cần thêm lời dạy, chúng cần bớt lời chì chiết.
\item Nhà là nơi về, đừng biến nó thành phòng thi kéo dài.
\item Cha mẹ sợ con khổ nên siết, nhưng siết quá đôi khi lại tạo ra một kiểu khổ khác.
\item Người lớn càng nóng càng nên học nói chậm với con.
\item Một đứa trẻ được lắng nghe đúng cách thường bớt chống đối hơn nhiều.
\item Gia đình hạnh phúc không phải nhà không có áp lực, mà là áp lực không giết mất tình người.
\end{enumerate}

\section{25 câu về biết ơn nhưng không tự xóa mình}
\begin{enumerate}[leftmargin=*,resume]
\item Biết ơn là cần, nhưng biết ơn không có nghĩa là em không được đau.
\item Đừng dùng chữ hiếu để buộc một người trẻ im luôn nhu cầu chính đáng của mình.
\item Mang ơn không đồng nghĩa với phải sống cả đời theo kịch bản người khác viết.
\item Có thể hiểu nỗi khổ của cha mẹ mà vẫn thừa nhận vết thương của mình.
\item Trưởng thành là vừa biết ơn vừa biết tự chịu trách nhiệm cho cuộc đời mình.
\item Nếu mọi lựa chọn của em chỉ để trả nợ cảm xúc cho gia đình, em sẽ rất khó sống thật.
\item Không ai cần con cái trả ơn bằng một cuộc đời sống trong oán âm thầm.
\item Yêu gia đình không phải là bỏ luôn quyền được khác.
\item Em có thể thương người nhà rất nhiều và vẫn cần khoảng cách lành mạnh.
\item Tự xóa mình để làm vừa lòng tất cả là một kiểu bất hiếu với chính đời mình.
\item Có người cả đời ngoan để rồi đến lúc nhìn lại chẳng biết mình thật sự muốn gì.
\item Biết ơn đúng là biến những gì nhận được thành sức sống tốt đẹp, không phải tự giam mình.
\item Đừng vì sợ mang tiếng mà chọn sống một phiên bản không phải mình.
\item Có những cuộc nói chuyện khó trong gia đình phải diễn ra, dù trễ.
\item Yêu thương trưởng thành luôn có phần rõ ràng chứ không chỉ chịu đựng.
\item Em không nợ ai một cuộc đời giả tạo.
\item Hiếu không nằm ở việc luôn nói dạ, mà còn ở việc sống đàng hoàng và không phá đời mình.
\item Tự lập cũng là một cách báo hiếu rất thực tế.
\item Nhiều người trẻ cần học cách nói thật với gia đình mà không bùng nổ.
\item Giữ hòa khí bằng cách nuốt hết mình vào trong không phải cách bền.
\item Có ranh giới với người nhà không làm em vô ơn.
\item Đôi khi điều gia đình cần không phải một đứa con hoàn hảo, mà là một đứa con rõ ràng hơn.
\item Chọn khác với kỳ vọng của nhà chưa chắc là phản bội.
\item Hòa thuận đẹp nhất là khi không ai phải mất mình quá nhiều.
\item Muốn thương lâu thì phải biết giữ mình nữa.
\end{enumerate}

\section{25 câu về vết thương trong nhà}
\begin{enumerate}[leftmargin=*,resume]
\item Có những câu nói trong nhà ngắn thôi mà theo người ta mười năm.
\item Trẻ con lớn nhanh trong tiếng quát, nhưng phần nào đó bên trong nó co lại.
\item Người lớn hay quên lời mình nói, con nít thường nhớ bằng cảm giác rất lâu.
\item Nhà mà ai cũng đúng quá thì chẳng ai dám thật.
\item Nhiều người trưởng thành phải dành nửa đời để chữa cái cách mình từng bị nuôi.
\item Tổn thương trong gia đình khó gọi tên vì nó đi kèm tình thương thật.
\item Có những đứa trẻ không cần hoàn hảo, chúng chỉ cần bớt bị làm cho thấy mình là gánh nặng.
\item Vết thương kín thường không ồn, nhưng nó lái hành vi rất mạnh.
\item Người hay tự trách quá mức thường từng bị chê quá nhiều ở nhà.
\item Cái lạnh trong nhà đáng sợ hơn cả tiếng ồn ngoài đường.
\item Không phải mọi nỗi đau đều cần đổ lỗi, nhưng mọi nỗi đau đều đáng được nhìn nhận.
\item Chữa lành không bắt đầu bằng tha thứ vội, mà bằng gọi đúng tên điều đã xảy ra.
\item Có người cả đời xin lỗi vì những thứ vốn không phải lỗi của họ.
\item Một cái ôm đúng lúc đôi khi sửa được rất nhiều thứ mà lời dạy không sửa nổi.
\item Tổn thương bị xem nhẹ thường lớn thêm trong im lặng.
\item Người bị làm cho thấy mình không đủ thường lớn lên với cơn đói công nhận rất mệt.
\item Không phải gia đình nào cũng biết yêu theo cách lành, nên người trẻ càng cần học cách tự lành.
\item Đừng chế giễu cảm xúc của con nít, nhiều vết nứt bắt đầu từ đó.
\item Ai cũng có thể thành phiên bản tệ của cha mẹ nếu không chịu nhìn lại những gì mình đã nhận.
\item Nhiều thói quen nóng nảy là di sản chưa được xử lý.
\item Một gia đình biết xin lỗi nhau đã cứu được rất nhiều đứa trẻ khỏi lớn lệch.
\item Nhà không cần hoàn hảo, chỉ cần có chỗ cho sự thật và sửa sai.
\item Nếu chưa được nghe lời tử tế, em càng phải học nói tử tế với mình.
\item Không ai chọn nơi mình sinh ra, nhưng có thể chọn cách không nhân bản mãi nỗi đau cũ.
\item Chữa lành là món quà em tặng cho cả mình và thế hệ sau.
\end{enumerate}

\section{25 câu về trưởng thành khỏi vòng lặp gia đình}
\begin{enumerate}[leftmargin=*,resume]
\item Trưởng thành là biết nhà mình có gì đẹp để giữ, có gì độc để cắt.
\item Không phải vì cha mẹ từng vậy mà em buộc phải sống y chang.
\item Hiểu nguồn gốc của vết thương không phải để nằm đó mãi.
\item Đừng lấy quá khứ đau làm bằng chứng rằng em không thể tử tế hơn.
\item Một người trẻ mạnh là người dừng được vài cái dở truyền đời.
\item Càng ý thức được vòng lặp cũ càng có cơ hội bẻ nó.
\item Đừng chỉ trách cách mình được nuôi, hãy học cách tự nuôi mình tốt hơn từ bây giờ.
\item Chọn khác đi đôi khi là một hành động rất yêu gia đình theo nghĩa dài hạn.
\item Muốn mai này nhà mình ấm hơn thì hôm nay phải bớt mang lửa cũ đi phát tiếp.
\item Không ai dạy em hết mọi kỹ năng sống, nhưng em có thể bắt đầu học muộn.
\item Tha thứ có thể đến sau, nhưng dừng cái sai nên đến sớm.
\item Nói câu này khó nhưng thật: có những điều người lớn đã làm sai với em.
\item Thừa nhận điều đó không làm em xấu, nó làm em tỉnh.
\item Đừng lấy lý do thương nhà mà chối bỏ nhu cầu chữa lành của mình.
\item Người trưởng thành không chỉ sống sót qua tuổi thơ, họ còn hiểu nó.
\item Mọi sự đổi khác đều bắt đầu từ một người dám không lặp lại.
\item Hôm nay em dịu hơn một chút đã là một bước khác quá khứ rồi.
\item Có những chiến thắng không ai thấy, như việc em không nói lại câu từng làm mình đau.
\item Bản lĩnh không phải cứng như nhà cũ dạy, mà là biết mềm đúng chỗ mà không gãy.
\item Đừng đợi đủ lành mới sống tiếp, hãy vừa sống vừa học cách lành.
\item Trưởng thành khỏi gia đình cũ không có nghĩa bỏ gia đình, mà là bớt để cái cũ điều khiển mình.
\item Em có quyền xây một ngôi nhà nội tâm tử tế hơn nơi mình từng nhận.
\item Quá khứ giải thích em, nhưng không được quyền viết hết tương lai em.
\item Mỗi phản ứng bớt độc là một viên gạch của đời mới.
\item Người dừng được vòng lặp đau chính là người âm thầm thay đổi cả một dòng sau mình.
\end{enumerate}

\chapter{100 câu cục súc về tiền bạc, thực tế và giá trị bản thân}
\markboth{100 câu cục súc về tiền bạc, thực tế và giá trị bản thân}{100 câu cục súc về tiền bạc, thực tế và giá trị bản thân}

\section{25 câu về tiền và sự tỉnh táo}
\begin{enumerate}[leftmargin=*]
\item Tiền không mua được tất cả, nhưng thiếu tiền làm nhiều thứ khó thấy mẹ.
\item Đừng chê tiền khi mọi giấc mơ của em đều cần tiền nuôi ít nhiều.
\item Yêu tiền không xấu, bán rẻ mình vì tiền mới đáng sợ.
\item Nghèo không làm ai cao quý tự động, giàu cũng không tự cho ai chiều sâu.
\item Tiền là công cụ, nhưng người rỗng rất dễ để nó thành ông chủ.
\item Kiếm tiền nhanh không sai, miễn là đừng mất luôn nhân cách và tương lai vì cái nhanh đó.
\item Đừng lấy câu sống một lần để tiêu như không có ngày mai.
\item Tiền vào tay người không biết giá trị sẽ chạy ra rất nhanh.
\item Nhiều người muốn tự do tài chính nhưng chưa tự do nổi trước ham muốn nhỏ.
\item Cái nghèo vật chất đã mệt, nghèo tư duy còn mệt dài.
\item Không phải thứ gì kiếm được nhiều cũng đáng để em dành tuổi trẻ cho.
\item Muốn có tiền sạch phải chịu đi đường không quá tắt.
\item Càng khó khăn càng phải học phân biệt cái cần và cái thèm.
\item Đừng tiêu để chứng minh mình ổn khi trong người đang rỗng ruột.
\item Nhiều món mua để đỡ buồn chỉ làm ví buồn thêm.
\item Biết quý đồng tiền khác hẳn sống khúm núm trước đồng tiền.
\item Tiền không nuôi nổi tự trọng nếu cách em kiếm hoặc tiêu đang làm em coi thường mình.
\item Đừng khoe tiêu mạnh nếu thu nhập của em đứng trên đôi chân rất yếu.
\item Cái giàu bền luôn chậm hơn cái giàu khoe.
\item Có những người muốn đời sang nhưng không muốn học kỹ năng xứng với mức sống đó.
\item Sống thực tế không phải giết mơ, là biết mơ cần hóa đơn nào phải trả.
\item Nỗi lo tiền bạc không biến mất bằng vài câu chữa lành nghe thơ.
\item Muốn bớt sợ tiền thì phải bớt mù về tiền.
\item Tiền rất thật, nhưng đừng để nó thành cái duy nhất làm em thấy mình có giá.
\item Cầm tiền cho đúng là một phần của trưởng thành, không phải chuyện người già.
\end{enumerate}

\section{25 câu về nghèo, xuất phát điểm và ý chí}
\begin{enumerate}[leftmargin=*,resume]
\item Xuất phát điểm thấp là bất lợi, không phải giấy miễn trách nhiệm.
\item Nghèo khiến đường dốc hơn, nhưng không tự động khiến em sâu sắc hay chăm hơn.
\item Đừng lãng mạn hóa cái nghèo, nó mệt thật.
\item Cũng đừng thần thánh hóa con nhà giàu, nhiều đứa có điều kiện nhưng đầu vẫn rỗng.
\item Khó khăn không phải cuộc thi xem ai than giỏi hơn.
\item Nhà không khá thì càng phải tỉnh trong lựa chọn, không phải càng được phép buông.
\item Có những người nghèo tiền nhưng giàu tự trọng, và đó là vốn rất lớn.
\item Xuất phát sau không có nghĩa là phải chạy loạn.
\item Càng ít chỗ dựa càng nên bớt mơ đường tắt.
\item Nghị lực không thay được hết điều kiện, nhưng thiếu nghị lực thì điều kiện tốt cũng uổng.
\item Đừng đố kỵ với người có nền, hãy học cách xây nền của mình từ chỗ đang đứng.
\item Hoàn cảnh giải thích nỗi khó, không nên được dùng để hợp thức hóa sự phó mặc mãi.
\item Người từ chỗ thấp đi lên thường không được phép sống mơ hồ quá lâu.
\item Bền hơn quan trọng hơn bốc lên một lúc rồi tắt.
\item Càng ít vốn càng phải đầu tư kỹ vào bản thân chứ đừng chỉ mơ trúng số đời.
\item Khổ không xấu, nhưng lấy khổ làm bản sắc rồi không chịu lớn thì phí.
\item Nhà nghèo mà còn tiêu kiểu sĩ diện thì tự trói thêm đá vào chân.
\item Đừng xấu hổ vì thiếu tiền, hãy xấu hổ nếu em thiếu cố gắng mà cứ trách đời.
\item Xuất phát điểm không do em chọn, nhưng cách đi tiếp thì có phần do em.
\item Càng thiếu ô dù càng cần tay nghề.
\item Có những người không thắng nhanh vì không có đường tắt, nhưng họ đi rất bền.
\item Đừng để cái nghèo dạy em thành người cay cú với mọi ai khá hơn.
\item Khó khăn nên làm em tỉnh, không nên làm em nát.
\item Thiếu tiền mà còn thiếu học thì gánh nặng tăng gấp đôi.
\item Người đi lên tử tế thường không rảnh khoe khổ, họ rảnh lo bước tiếp hơn.
\end{enumerate}

\section{25 câu về giá trị bản thân không bán rẻ}
\begin{enumerate}[leftmargin=*,resume]
\item Đừng định giá mình bằng mức người khác sẵn sàng lợi dụng.
\item Cần tiền không có nghĩa là nhận mọi cuộc chơi bẩn.
\item Cái rẻ nhất một người trẻ hay bán là lòng tự trọng trong vài phút yếu lòng.
\item Giá trị bản thân không nằm ở việc em được bao nhiêu người cần, mà ở việc em giữ được gì khi bị cần.
\item Có giá trị thì nói ít thôi, làm cho chắc.
\item Đừng cúi quá thấp trước người có thứ em đang thiếu.
\item Thiếu tiền có thể sửa, mất phẩm giá sửa khó hơn nhiều.
\item Người biết mình đáng gì sẽ bớt vội nhận mọi lời mời gọi hào nhoáng.
\item Muốn người khác tôn trọng thì đừng tự xài mình như đồ rẻ tiền.
\item Có nguyên tắc nghe bất tiện lúc đầu nhưng cứu em khỏi rất nhiều lần tự khinh mình về sau.
\item Bán sức lao động là bình thường, bán luôn cốt cách thì khác.
\item Giá trị thật ít khi ồn ào, nhưng rất khó thay thế.
\item Người trẻ hay lầm giữa được chú ý và được tôn trọng.
\item Đừng vì một khoản trước mắt mà ký vào cuộc đời dài hạn của mình quá dễ.
\item Chỗ nào bắt em chà đạp mình để được vào, chỗ đó đáng nghi.
\item Giá trị bản thân không phải câu quote treo tường, nó là cách em ra quyết định lúc túng.
\item Khi cùng quẫn, nhân cách mới lộ rõ mình luyện được tới đâu.
\item Không ai giữ giá cho em nếu chính em liên tục tự giảm giá mình.
\item Cần được công nhận là thật, nhưng đừng bán mình chỉ để đổi lấy tiếng wow ngắn hạn.
\item Tử tế với mình đi rồi người khác mới khó coi thường em hơn.
\item Đứng thẳng không làm em giàu ngay, nhưng giữ được đường dài.
\item Có vài lời từ chối nghe mất tiền nhưng thực ra giữ vốn sống.
\item Sự đàng hoàng của em có thể không được trả công ngay, nhưng nó xây uy tín âm thầm.
\item Giá trị càng thật càng không cần trang trí nhiều.
\item Đừng biến bản thân thành món rẻ trong phiên chợ ngắn hạn của người khác.
\end{enumerate}

\section{25 câu về sống thực tế mà không mất chiều sâu}
\begin{enumerate}[leftmargin=*,resume]
\item Sống thực tế không phải là chê bai mọi điều đẹp.
\item Thực tế là biết cái nào mơ được, cái nào phải trả giá.
\item Đừng gọi mình thực tế nếu em chỉ đang sợ dấn thân.
\item Có người bảo thẳng tính, thực ra là họ mất niềm tin nên cái gì cũng chua.
\item Thực tế tử tế là nhìn đời rõ mà không thành người cạn.
\item Người càng hiểu thực tế càng ít nói mấy câu tuyệt đối.
\item Biết đời khó không có nghĩa là ngừng cố sống đẹp.
\item Thực tế không bắt em vô cảm, nó chỉ bắt em bớt ảo.
\item Đừng để câu ngoài đời là thế giết mất tiêu chuẩn sống đàng hoàng trong em.
\item Mơ cao mà không tính tiền là hão, chỉ biết tiền mà không biết mình muốn thành ai cũng héo.
\item Người sâu là người nhìn thẳng vào hóa đơn của cuộc đời mà vẫn giữ được phần người.
\item Sống thực tế mà không biết yêu gì ngoài lợi ích thì rất nhanh khô.
\item Mỗi người đều cần một cái la bàn ngoài chuyện kiếm sống.
\item Chín chắn không phải lạnh, là biết thứ gì đáng đặt tim và sức vào.
\item Có nguyên tắc giữa đời thực dụng là một kiểu giàu ít người thấy.
\item Đừng vì gặp nhiều thứ dơ mà tập quen với việc sống dơ.
\item Thực tế tốt sẽ làm em bền, thực tế độc chỉ làm em mỏi.
\item Đừng chế giễu người còn tin vào điều đẹp nếu họ cũng biết lao động cho điều đó.
\item Vừa tỉnh vừa ấm là khó, nhưng đáng luyện.
\item Người biết đời không đơn giản thường cũng bớt phán nhanh người khác.
\item Thực tế mà không mất nhân tính mới là trình độ.
\item Bước chân xuống đất nhưng đừng để mắt chỉ nhìn bùn.
\item Có tiền là một chuyện, sống có giá trị là chuyện khác nhưng không nên tách rời nhau quá xa.
\item Cứng cỏi không đồng nghĩa với mất khả năng rung động.
\item Sống thật với đời không cần phải sống cộc với mọi điều tử tế.
\end{enumerate}